% preamble/part.title.tex
% Part 扉页设计 —— 墨绿 + 暗金
% Arnold 的常微分方程原书未分 Part，此文件保留以备扩展使用。
% =========================================================

\RequirePackage{tikz}
\RequirePackage{graphicx}
\RequirePackage{etoolbox}
\RequirePackage{setspace}

\makeatletter

\gdef\partimage{}
\newcommand{\setpartimage}[1]{\gdef\partimage{#1}}

\gdef\partintro{}
\newcommand{\setpartintro}[1]{\gdef\partintro{#1}}

\newlength{\PartRibbonHeight}
\setlength{\PartRibbonHeight}{2.35cm}

\newlength{\PartContentTop}
\setlength{\PartContentTop}{\dimexpr\PartRibbonHeight+1.2cm\relax}

\newlength{\PartDividerY}
\setlength{\PartDividerY}{\dimexpr\PartRibbonHeight+6.8cm\relax}

\newcommand{\PartTitleBlock}[2]{%
    \ifdefempty{#1}{%
        \fill[GoldMain!55]
            ([xshift=1.8cm,yshift=-\PartContentTop]current page.north west)
            -- ++(0.22,0.22) -- ++(-0.22,0.22) -- ++(-0.22,-0.22) -- cycle;
    }{%
        \node[anchor=north west,text=GoldMain!72,font=\fontsize{52}{62}\bfseries]
            at ([xshift=1.8cm,yshift=-\PartContentTop]current page.north west)
            {#1};
        \node[anchor=north west,text=GoldMain!48,font=\Large\bfseries]
            at ([xshift=1.8cm,yshift=-\PartContentTop-3.1cm]current page.north west)
            {Part};
    }%
    \node[anchor=north west,text=BlueDark,font=\fontsize{22}{28}\bfseries,
          text width=8cm,align=left]
        at ([xshift=1.8cm,yshift=-\PartContentTop-4.4cm]current page.north west)
        {#2};
}

\newcommand{\PartIntroBlock}{%
    \ifdefempty{\partintro}{}{%
        \node[anchor=north east,text=BlueMain!80,font=\Large\bfseries]
            at ([xshift=-1.8cm,yshift=-\PartContentTop]current page.north east)
            {简\quad 介};
        \node[anchor=north east,text=BlueDark!55,font=\large\itshape,
              text width=7.4cm,align=right,inner sep=0pt]
            at ([xshift=-1.8cm,yshift=-\PartContentTop-1.5cm]current page.north east)
            {\setstretch{1.45}\partintro};
    }%
}

\newcommand{\PartTitlePage}[2]{%
    \thispagestyle{empty}%
    \vspace*{-\baselineskip}%
    \begin{tikzpicture}[remember picture,overlay]
        \fill[BlueMain!94]
            (current page.north west) rectangle
            ([yshift=-\PartRibbonHeight]current page.north east);
        \fill[GoldMain!75]
            (current page.north west) rectangle
            ([xshift=0.18cm,yshift=-\PartRibbonHeight]current page.north west);
        \fill[GoldMain!75]
            ([xshift=-0.18cm]current page.north east) rectangle
            ([yshift=-\PartRibbonHeight]current page.north east);
        \PartTitleBlock{#1}{#2}
        \PartIntroBlock
        \draw[GoldMain!20,line width=0.5pt,dash pattern=on 3pt off 5pt]
            ([xshift=11.0cm,yshift=-\PartContentTop]current page.north west) --
            ([xshift=11.0cm,yshift=-\PartDividerY+0.3cm]current page.north west);
        \draw[GoldMain!42,line width=0.7pt]
            ([yshift=-\PartDividerY]current page.north west) --
            ([yshift=-\PartDividerY]current page.north east);
        \fill[GoldMain!50]
            ([yshift=-\PartDividerY]current page.north west) circle (1.8pt);
        \fill[GoldMain!50]
            ([yshift=-\PartDividerY]current page.north east) circle (1.8pt);
        \fill[GoldMain!60]
            ([yshift=-\PartDividerY]current page.north)
            -- ++(0.13,0.13) -- ++(-0.13,0.13) -- ++(-0.13,-0.13) -- cycle;
        \ifdefempty{\partimage}{%
            \shade[top color=BlueMain!12,bottom color=BlueMain!3]
                ([yshift=-\PartDividerY]current page.north west) rectangle
                (current page.south east);
            \node[align=center,text=BlueDark!40,font=\large\itshape]
                at ([yshift=-5.5cm]current page.center)
                {请在 \texttt{\textbackslash part} 前设置图片};
        }{%
            \node[anchor=north west,inner sep=0pt,outer sep=0pt]
                at ([yshift=-\PartDividerY]current page.north west)
                {\includegraphics[
                    width=\paperwidth,
                    height=\dimexpr\paperheight-\PartDividerY\relax
                ]{\partimage}};
        }%
    \end{tikzpicture}%
    \clearpage%
    \if@twoside
        \if@openright
            \null
            \thispagestyle{empty}%
            \newpage
        \fi
    \fi
}

\def\@part[#1]#2{%
    \ifnum \c@secnumdepth >-2\relax
        \refstepcounter{part}%
        \addcontentsline{toc}{part}{\thepart\hspace{1em}#1}%
    \else
        \addcontentsline{toc}{part}{#1}%
    \fi
    \markboth{}{}%
    \PartTitlePage{\thepart}{#2}%
    \@endpart
}

\def\@spart#1{%
    \PartTitlePage{}{#1}%
    \@endpart
}

\def\@endpart{%
    \if@tempswa
        \twocolumn
    \fi
}

\makeatother
